%% Chapitre 05 — écrit par un subagent Rédacteur.
%% RÈGLE : uniquement des commandes sémantiques bookforge. Aucun \chapter,
%% \textbf, \emph, \begin{lstlisting}, \vspace... brut.

\chapitre{Gérer le contexte : la ressource qui décide de tout}

\introChapitre{
  Un agent ne raisonne bien que sur ce qu'il a sous les yeux. Trop peu de
  contexte et il invente ; trop et il se noie ou coûte cher. Ce chapitre
  durable explique comment l'agent acquiert, retient et perd du contexte, et
  comment tu peux le piloter activement plutôt que le subir.
}

\section{Ce que l'agent garde en tête pendant une session}

\paragraphe{Au chapitre précédent, tu as appris à découper ton travail en
tâches délégables. Mais une tâche bien formulée ne suffit pas : encore
faut-il que l'agent ait, au bon moment, les bons éléments sous les yeux. C'est
exactement le problème du contexte, et c'est là que beaucoup de délégations
ratent. Tu donnes un but clair, puis l'agent répond à côté, non parce qu'il est
bête, mais parce qu'il ne voit pas ce que toi tu vois.}

\paragraphe{Pendant une session, le contexte de l'agent est un empilement.
Il contient tes instructions, l'historique de l'échange, les fichiers qu'il a
lus, les résultats des commandes qu'il a lancées. Tout cela s'accumule au fil
des messages. L'agent ne raisonne que là-dessus : ce qui n'y figure pas
n'existe pas pour lui. Un fichier que tu n'as pas mentionné et qu'il n'a pas
ouvert ? Inconnu. Une décision prise oralement il y a deux jours ? Disparue.}

\paragraphe{Retiens ce principe : \important{le contexte est une fenêtre, pas
une mémoire totale.} L'agent ne connaît pas ton projet « en général ». Il
connaît le morceau de projet qu'il a chargé dans cette session-ci. Ton rôle,
c'est de décider quel morceau.}

\section{Comment l'agent va chercher du contexte}

\paragraphe{Bonne nouvelle : tu n'as pas à tout lui copier-coller. Un agent
comme Claude Code va chercher lui-même ce dont il a besoin. Il liste
les fichiers, lit ceux qui semblent pertinents, recherche un symbole ou une
chaîne dans le dépôt, lance une commande pour observer un comportement. C'est
toute la différence avec un simple chatbot : il agit pour s'informer avant
d'agir pour produire.}

\paragraphe{Mais cette autonomie a une contrepartie. L'agent explore à partir
des indices que tu lui donnes. Si tu écris « corrige le bug d'authentification »
sans plus, il devine où chercher, et il devine parfois mal : il ouvre dix
fichiers, en lit trois de travers, et reconstitue une image approximative du
système. Tu paies cette exploration en lenteur, en coût, et en risque
d'erreur.}

\begin{astuce}
Amorce l'exploration au lieu de la laisser tâtonner. Nomme les fichiers ou
dossiers concernés, donne le nom exact de la fonction, colle le message
d'erreur complet. Tu ne fais pas le travail à sa place : tu lui évites de
fouiller à l'aveugle, et tu obtiens une réponse plus juste, plus vite.
\end{astuce}

\paragraphe{Le bon réflexe n'est donc ni de tout fournir, ni de tout laisser
deviner. C'est de \emphase{pointer} : indiquer le point d'entrée, puis laisser
l'agent suivre le fil depuis là.}

\section{La fenêtre de contexte : limites, saturation, dégradation}

\paragraphe{Cet espace de travail a une taille maximale. On l'appelle la
\terme{fenêtre de contexte} : la quantité d'information que l'agent peut garder
active en même temps. Elle est grande, mais finie. Et au cours d'une longue
session, on la remplit plus vite qu'on ne croit.}

\paragraphe{Chaque fichier lu, chaque sortie de commande verbeuse, chaque
aller-retour s'y entasse. Quand la fenêtre approche de la saturation, deux
choses arrivent. D'abord, l'agent commence à \emphase{oublier} : les éléments
les plus anciens de la conversation sortent de sa vue. La consigne précise que
tu avais donnée au début ? Il ne l'a peut-être plus. Ensuite, même sans
oubli franc, sa qualité se \emphase{dégrade} : noyé sous trop d'information, il
peine à distinguer l'essentiel du bruit et ses réponses deviennent plus floues.}

\begin{avertissement}
Le symptôme typique d'une fenêtre saturée : l'agent répète une erreur que tu
viens de corriger, redemande une information déjà donnée, ou contredit une
décision prise plus tôt dans la même session. Ce n'est pas de la mauvaise
volonté, c'est un débordement. Insister ne sert à rien : il faut alléger le
contexte.
\end{avertissement}

\paragraphe{Comprendre cette limite change ta manière de travailler. Une
session n'est pas un sac sans fond où entasser tout un projet. C'est un plan de
travail de surface limitée. \important{Plus tu gardes ce plan dégagé, mieux
l'agent travaille.}}

\section{La mémoire projet : des instructions qui survivent à la session}

\paragraphe{Certaines informations ne devraient jamais être réexpliquées à
chaque session : tes conventions de code, les commandes pour tester, les
pièges connus du dépôt, ce qu'il ne faut surtout pas toucher. Les répéter à la
main serait absurde. C'est le rôle de la \terme{mémoire projet} : un fichier
d'instructions persistantes, rangé dans le dépôt, que l'agent charge
automatiquement au démarrage de chaque session sur ce projet.}

\paragraphe{Tu y écris une fois ce qui vaut pour toujours. Par exemple :}

\begin{extraitcode}{text}
# Conventions du projet

- Lance les tests avec `make test` avant de proposer un diff.
- N'edite jamais les fichiers du dossier `generated/` : ils sont produits
  automatiquement.
- Style : fonctions courtes, pas de dependance nouvelle sans la justifier.
- La logique metier vit dans `src/core/`, jamais dans les handlers HTTP.
\end{extraitcode}

\paragraphe{Cette mémoire devient le socle de contexte commun à toutes tes
délégations. Tu n'as plus à rappeler « teste avant de me montrer le résultat » :
c'est acquis. Mais garde-la \emphase{courte et vraie}. Une mémoire projet
gonflée de détails périmés est pire que pas de mémoire : elle occupe la fenêtre
et désinforme l'agent.}

\begin{astuce}
Traite la mémoire projet comme du code : relis-la, élague ce qui ne sert plus,
corrige ce qui a changé. Une bonne règle : si une instruction t'a fait gagner
du temps trois fois, elle mérite d'y figurer ; si tu ne la suis plus toi-même,
elle doit en sortir.
\end{astuce}

\section{Repartir propre : réinitialiser le contexte}

\paragraphe{Tôt ou tard, la session devient encombrée. Tu as exploré trois
pistes, abandonné deux, fait des essais sans suite. Tout ce passé traîne dans
le contexte et brouille l'agent. La solution n'est pas de continuer à pousser :
c'est de \important{repartir propre}.}

\paragraphe{Réinitialiser, c'est ouvrir une session neuve, vidée de l'historique
accumulé. L'agent recharge la mémoire projet, mais oublie le fouillis de
l'échange précédent. C'est libérateur : tu repars d'un plan de travail dégagé,
avec une formulation enfin claire de ce que tu veux, débarrassée des fausses
pistes.}

\paragraphe{Le réflexe utile : quand tu changes de tâche, change de session.
Garder une seule conversation interminable pour cinq sujets différents, c'est
la garantie d'une fenêtre saturée et d'un agent confus. Une tâche délégable,
une session propre. Avant de réinitialiser, capture juste ce qui mérite de
survivre : dans la mémoire projet, dans une note, ou dans un diff déjà
appliqué.}

\section{Les signes qu'il faut recadrer ou redémarrer}

\paragraphe{Apprends à reconnaître quand le contexte travaille contre toi. Voici
les signaux qui doivent te faire réagir :}

\begin{liste}
  \element{L'agent répète une erreur que tu viens de corriger, ou redemande une
  information déjà fournie : la fenêtre déborde.}
  \element{Ses réponses deviennent vagues, génériques, comme s'il avait perdu le
  fil de l'objectif : il est noyé.}
  \element{Il s'accroche obstinément à une approche que tu as écartée : le passé
  de la session le tire en arrière.}
  \element{Tu passes plus de temps à le recadrer qu'à avancer : la session est
  usée.}
\end{liste}

\paragraphe{Face à ces signaux, deux gestes. Le plus léger : \emphase{recadrer}
sans tout jeter, en reformulant l'objectif et en repointant les bons fichiers.
Le plus net : \emphase{redémarrer} une session propre quand le fouillis est trop
épais. Ni l'un ni l'autre n'est un aveu d'échec ; piloter le contexte fait
partie du métier de déléguer.}

\resumeChapitre{
  \element{\important{Fenêtre, pas mémoire} : l'agent ne connaît que ce qu'il a chargé dans la session en cours — tout le reste n'existe pas pour lui.}
  \element{\important{Pointer, pas copier} : indique un chemin de fichier plutôt que de coller son contenu — l'agent lira la version réelle et courante.}
  \element{\important{Saturation} : quand la fenêtre déborde, l'agent oublie les décisions anciennes et ses réponses se dégradent ; le symptôme est de répéter une erreur déjà corrigée.}
  \element{\important{Mémoire projet} : écris dans un fichier versionné les conventions et contraintes durables du dépôt — l'agent les charge automatiquement à chaque session.}
  \element{\important{Repartir propre} : change de session quand tu changes de tâche ; une tâche, une session — c'est le réflexe qui préserve la qualité du contexte.}
}

\paragraphe{Tu sais désormais donner à l'agent ce qu'il faut voir, ni plus ni
moins, et reconnaître quand il faut faire le ménage. Reste la question qui
décide de tout : comment savoir si ce qu'il te rend est juste ? Car même avec un
contexte parfait, l'agent peut se tromper avec aplomb. Le chapitre suivant
s'attaque à la vérification de son travail.}
